\chapter{三相记忆的深度理解：从存储分类到认知结构持续化动力学}

\begin{chapterlead}
前一章已经从信息相态的角度区分了三相记忆 $h$、$E$ 与 $\Theta$，回答了它们在信息组织方式上“有何不同”。本章在第二编的理论脉络中继续向更深处推进，追问一个更根本的问题：三类记忆的真正分野，并不首先在于存储位置或时间长短，而在于过去以何种动力深度继续存在于未来。为此，我们把记忆的定义从“把信息保存在某个位置”改写为“过去事件对未来认知动力保持因果影响的能力”，把研究问题从“记忆存在哪里”转向“过去如何继续存在”。由此，三相记忆不再是一种存储分类，而成为一套认知结构持续化动力学：三相不是三个互不相关的容器，而是同一认知结构在时间中逐渐获得不同持续方式的三个主要动力相区。

本章的论证分三步展开。第一步逐相重建本体刻画：$h$ 的本质是活动性而非固定时长——只要一个结构仍直接参与当前因果闭环，它在功能上仍然属于现在；$E$ 是已经退出当前活动、但仍保持事件个体性并可重新识别的过去，其本质是关系实例而非内容复制，因而处于既保持历史身份又允许未来重解释的“亚稳”液态；$\Theta$ 是由历史沉降形成、已经成为未来生成条件的慢结构，它同时属于过去与未来。第二步把三相放回同一时间尺度谱：核心过程不是“保存、搬运、长期保存”，而是从活动到经历身份再到生成结构的相变，并且这一过程可以反向进行——召回与巩固是相反方向的时间尺度迁移，$\Theta\rightarrow E$ 还包括预测与想象，三相由此构成再生产循环 $\Theta_t\rightarrow h_t\rightarrow E_t\rightarrow\Theta_{t+1}$，而不是单向层级。第三步使上述刻画可度量：我们定义当前激活度、情景身份保持度与生成约束强度三个序参量，并说明工作记忆的有限性为什么会自然要求三相结构、情景层为什么是快速经历与慢结构之间的保护层、遗忘为什么不是单一现象。

本章的结论是：三相记忆的最深层区别不是“短期、中期、长期”，而是现在、历史身份与生成规律——$h$ 是仍然作为现在存在的过去，$E$ 是仍然能够被识别的过去，$\Theta$ 是已经成为未来生成结构的过去；记忆不是保存过去，而是让过去以不同速度继续存在于未来。这一动力学理解同时为后文做好准备：当三相结构随时间不断推进并循环再生产，记忆体就不再是沿物理时间追加的档案，而将显现为一个具有内部结构时间的生命周期整体，这正是下一章要建立的生命周期记忆体理论的出发点。
\end{chapterlead}

\section{从“记忆存在哪里”转向“过去如何继续存在”}

在三相记忆理论的早期描述中，工作记忆 $h$、情景记忆 $E$ 与结构记忆 $\Theta$ 很容易被理解为三类不同的存储系统：
\[
h=\text{短期工作空间},
\qquad
E=\text{情景存储},
\qquad
\Theta=\text{长期结构存储}.
\]

这种理解在工程上是直观的，但在理论上仍然不够深入。它默认了一个传统存储观念，即记忆的本质是“把信息保存在某个位置”。然而，对于持续智能体而言，更根本的问题并不是
\[
\text{过去被存在哪里？}
\]
而是
\begin{statementbox}
\centering
\adjustbox{max width=.96\linewidth}{\(\displaystyle
\text{过去发生过的事情，如何继续改变未来的状态演化？}
\)}
\end{statementbox}

因此，记忆可以首先被重新定义为：
\begin{statementbox}
\centering
\adjustbox{max width=.96\linewidth}{\(\displaystyle
\text{过去事件对未来认知动力保持因果影响的能力。}
\)}
\end{statementbox}

只要一个过去状态仍然能够影响未来状态转移，它就以某种形式继续存在于系统中。由此，三相记忆不再首先是一种存储层级，而是一种描述“过去以何种深度进入未来”的动力学理论。

三相所对应的，不是三个互不相关的容器，而是同一认知结构在时间中逐渐获得不同持续方式的三个主要动力区域：
\[
\boxed{
h
\rightarrow
E
\rightarrow
\Theta.
}
\]

其中，$h$ 表示过去尚未退出当前活动，$E$ 表示过去已经退出当前活动但仍保持事件个体性，而 $\Theta$ 表示过去已经沉降为塑造未来状态的生成结构。

\section{$h$：过去尚未离开现在}

工作记忆 $h$ 最容易被误解为一个“短期存储区”。更准确地说，$h$ 表示当前仍然直接参与认知因果闭环的结构。

因此：
\begin{statementbox}
\centering
\adjustbox{max width=.96\linewidth}{\(\displaystyle
h=
\text{仍在发生的认知结构}.
\)}
\end{statementbox}

它并不一定已经成为“记忆中的过去”。例如，一个当前正在被比较的对象、一个尚未完成的推理、一段仍然影响当前判断的感觉、一个当前被维持的目标，都属于 $h$。

从时间关系看，可以写成：
\begin{statementbox}
\centering
\adjustbox{max width=.96\linewidth}{\(\displaystyle
h \subset Past\cap Present.
\)}
\end{statementbox}

这里的“过去”是指这些状态可能在几秒、几百毫秒甚至更早已经开始，但只要其影响尚未退出当前认知闭环，它们仍然属于现在。

因此，判断一个结构是否属于 $h$，核心问题不是“它出现多久了”，而是：
\begin{statementbox}
\centering
\adjustbox{max width=.96\linewidth}{\(\displaystyle
\text{它是否仍然直接参与当前状态转移？}
\)}
\end{statementbox}

这比“短期”更准确。

例如，某个目标可能已经维持几十分钟，只要它仍然直接参与当前规划，它在功能上仍然属于工作记忆的活动面；反过来，一个仅持续数秒的刺激，一旦不再参与当前认知，也可以已经退出 $h$。

因此：
\begin{statementbox}
\centering
\adjustbox{max width=.96\linewidth}{\(\displaystyle
h\text{ 的本质是活动性，而不是固定时间长度。}
\)}
\end{statementbox}

\section{$E$：已经离开现在，但仍保持经历身份的过去}

当一个状态不再直接参与当前认知，它便退出了 $h$。但这并不意味着它必须消失。

如果系统仍然能够在未来重新识别：
\[
\text{某件事情曾经发生过},
\]
并能够恢复与该事件有关的时间、对象、关系、结果或上下文，那么过去已经形成了情景记忆 $E$。

因此：
\begin{statementbox}
\centering
\adjustbox{max width=.96\linewidth}{\(\displaystyle
E=
\text{已经退出当前活动，但仍保持事件个体性的过去结构}.
\)}
\end{statementbox}

可以进一步写成：
\begin{statementbox}
\centering
\adjustbox{max width=.96\linewidth}{\(\displaystyle
E=
\text{Re-identifiablePast}.
\)}
\end{statementbox}

这里“可重新识别”是关键。

例如，以下两种状态存在根本差异：

第一种：
\[
\text{某次活动在系统中留下了模糊残余影响};
\]

第二种：
\[
\text{系统能够重新识别“那一次发生了什么”}.
\]

只有第二种才构成完整意义上的情景记忆。

因此必须区分：
\begin{statementbox}
\centering
\adjustbox{max width=.96\linewidth}{\(\displaystyle
Trace\neq Episode.
\)}
\end{statementbox}

一个简单残留状态可以说明“过去留下过影响”，但情景记忆要求过去被组织为一个具有边界的关系实例。

\section{情景记忆的本质是关系实例，而不是内容复制}

对情景记忆更深入的理解是：它并不需要完整复制事件中所有对象本身的知识，而主要保存这些对象在某一次具体时间和环境中形成了怎样的关系。

设系统已经知道若干稳定认知结构：
\[
C_1,C_2,\ldots,C_n.
\]

一次具体经历可以形成：
\[
E_i=
\operatorname{Bind}
(
C_1,C_2,\ldots,C_n,
t_i,
R_i
),
\]
其中 $t_i$ 表示时间结构，$R_i$ 表示这一具体事件中的关系结构。

因此：
\begin{statementbox}
\centering
\adjustbox{max width=.96\linewidth}{\(\displaystyle
E_i=
\text{一次具体的、带时间索引的关系实例化}.
\)}
\end{statementbox}

例如，一个系统长期已经知道“杯子”“厨房”“某个人”“放置”这些概念。某一次经历并不需要重新存储这些概念本身，而是保存：
\[
\text{某个人在某个时间把某个杯子放在厨房}.
\]

所以：
\begin{statementbox}
\centering
\adjustbox{max width=.96\linewidth}{\(\displaystyle
E
\text{ 的核心不是对象本体，而是对象关系的历史实例。}
\)}
\end{statementbox}

这也解释了为什么情景记忆天然具有：
\[
\text{What},
\quad
\text{Where},
\quad
\text{When},
\quad
\text{Relation}
\]
等结构。

\section{$E$ 为什么处于“液态”}

将 $E$ 称为“液态”，并不仅仅是因为其时间尺度位于 $h$ 与 $\Theta$ 之间。

更深层的原因在于，情景记忆既保持某种结构身份，又保持高度可塑性。

一个 Episode 可以被：
\[
\text{重新解释},
\quad
\text{合并},
\quad
\text{拆分},
\quad
\text{弱化},
\quad
\text{遗忘},
\quad
\text{重新赋予重要性},
\]
而不会像当前工作状态那样瞬时变化，也不会像结构记忆那样成为未来认知的强约束。

因此，可以把 $E$ 定义为：
\begin{statementbox}
\centering
\adjustbox{max width=.96\linewidth}{\(\displaystyle
E=
\text{MetastableStructuredPast}.
\)}
\end{statementbox}

即：
\begin{statementbox}
\centering
\adjustbox{max width=.96\linewidth}{\(\displaystyle
\text{亚稳的、结构化的过去}.
\)}
\end{statementbox}

“亚稳”意味着它可以持续较长时间，但仍然允许重新组织；“结构化”意味着它已经不再只是模糊残留，而是能够被重新识别和调用的经历单位。

因此 $E$ 的重要特性不是“中期”，而是：
\begin{statementbox}
\centering
\adjustbox{max width=.96\linewidth}{\(\displaystyle
\text{既保持历史个体性，又允许未来重解释。}
\)}
\end{statementbox}

\section{$\Theta$：已经进入未来的过去}

结构记忆 $\Theta$ 最容易被误解为“长期知识存储”。

更准确地说：
\begin{statementbox}
\centering
\adjustbox{max width=.96\linewidth}{\(\displaystyle
\Theta
\text{ 不是过去的收藏，而是过去已经改变了系统未来如何演化。}
\)}
\end{statementbox}

因此：
\begin{statementbox}
\centering
\adjustbox{max width=.96\linewidth}{\(\displaystyle
\Theta=
\text{PastSedimentedIntoFutureDynamics}.
\)}
\end{statementbox}

中文可以表述为：
\begin{statementbox}
\centering
\adjustbox{max width=.96\linewidth}{\(\displaystyle
\text{过去沉降成为未来动力规律后的结构。}
\)}
\end{statementbox}

例如，一个智能体经过大量具体经验后，不再需要逐次回忆每一个事件，却已经形成了稳定的：
\[
\text{对象规律},
\quad
\text{因果关系},
\quad
\text{行为模式},
\quad
\text{技能},
\quad
\text{预测结构}.
\]

这些结构已经进入 $\Theta$。

因此 $E$ 和 $\Theta$ 的根本区别不是“保存时间更长”，而是：
\begin{statementbox}
\centering
\adjustbox{max width=.96\linewidth}{\(\displaystyle
E
\text{ 保持“那一次”的个体性};
\)}
\end{statementbox}
\begin{statementbox}
\centering
\adjustbox{max width=.96\linewidth}{\(\displaystyle
\Theta
\text{ 失去具体事件个体性，却获得更强的生成能力}.
\)}
\end{statementbox}

\section{从情景到结构：个体性下降，生成性上升}

设系统积累了一系列具体经历：
\[
E_1,E_2,\ldots,E_n.
\]

这些经历可能具有共同结构。经过长期巩固后，它们可以形成：
\[
\Theta_C.
\]

这一过程可以写成：
\[
\boxed{
E_1+E_2+\cdots+E_n
\rightarrow
\Theta_C.
}
\]

但这里并不是简单的信息相加。

在这一过程中：
\begin{statementbox}
\centering
\adjustbox{max width=.96\linewidth}{\(\displaystyle
\text{EpisodeIdentity}\downarrow
\)}
\end{statementbox}
同时：
\begin{statementbox}
\centering
\adjustbox{max width=.96\linewidth}{\(\displaystyle
\text{GenerativePower}\uparrow.
\)}
\end{statementbox}

也就是说，具体经历的细节和个体性逐渐被压缩，而其中稳定、重复、具有预测价值的结构被保留下来。

因此：
\begin{statementbox}
\centering
\adjustbox{max width=.96\linewidth}{\(\displaystyle
E\rightarrow\Theta
=
\text{Instance}\rightarrow\text{Generator}.
\)}
\end{statementbox}

这比“短期记忆转化为长期记忆”更加准确。

例如，一个智能体经历多次不同的门以后，逐渐形成“门把手通常需要旋转并随后推或拉”的结构知识。此时，它可能已经不记得每一次具体开门过程，但能力反而增强。

因此：
\begin{statementbox}
\centering
\adjustbox{max width=.96\linewidth}{\(\displaystyle
\text{遗忘具体经历}
\neq
\text{丧失学习成果}.
\)}
\end{statementbox}

甚至在许多情况下：
\[
\boxed{
E\downarrow
}
\]
恰恰可能意味着：
\[
\boxed{
E\rightarrow\Theta
}
\]
已经成功完成。

\section{记住一件事与学会一件事不是同一个过程}

三相记忆使“记忆”和“学习”获得了明确区分。

记住一次经历，主要表现为：
\[
\boxed{
E\rightarrow h.
}
\]

即过去的具体经历重新进入当前活动。

而学习则主要表现为：
\[
\boxed{
E\rightarrow\Theta.
}
\]

即经历中的稳定结构改变了系统未来的生成规律。

因此：
\begin{statementbox}
\centering
\adjustbox{max width=.96\linewidth}{\(\displaystyle
\text{Remembering}\neq\text{Learning}.
\)}
\end{statementbox}

一个系统可以记得很多经历，但没有从中形成稳定结构；也可以已经形成稳定能力，却无法回忆其形成过程中所有具体经历。

这种区分对于持续智能体尤其重要，因为长期智能不可能依靠无限增加可检索 Episode 来获得能力。真正成熟的系统必须不断把具体经历转化为更加压缩的生成结构。

\section{三相不是三个容器，而是时间尺度谱上的三个主要相区}

将 $h$、$E$、$\Theta$ 画成三个独立盒子虽然便于理解，但容易产生误导。

更一般的持续智能系统可能具有一系列连续的时间尺度：
\[
\tau_1<\tau_2<\cdots<\tau_n.
\]

一个认知结构可能逐步经历：
\[
\text{Active}
\rightarrow
\text{Persistent}
\rightarrow
\text{Metastable}
\rightarrow
\text{Consolidated}
\rightarrow
\text{Structural}.
\]

因此 $h$、$E$ 和 $\Theta$ 更适合被理解为这一连续时间尺度谱中的三个主要动力区域：
\begin{statementbox}
\centering
\adjustbox{max width=.96\linewidth}{\(\displaystyle
h,E,\Theta
=
\text{ThreeDominantMemoryPhaseRegions}.
\)}
\end{statementbox}

可以定义一个认知结构 $C$ 的有效时间常数：
\[
\tau(C,t).
\]

则：
\[
h:
\quad
\tau(C,t)\sim\tau_{\mathrm{fast}},
\]
\[
E:
\quad
\tau_{\mathrm{fast}}
\ll
\tau(C,t)
\ll
\tau_{\mathrm{slow}},
\]
\[
\Theta:
\quad
\tau(C,t)\sim\tau_{\mathrm{slow}}.
\]

这里的边界不是绝对固定的数值，而是不同智能体、不同实现和不同生命周期中相对稳定的时间尺度分离。

\section{从“记忆搬运”转向“记忆相变”}

如果三相是同一认知结构在不同时间尺度上的存在方式，那么
\[
h\rightarrow E\rightarrow\Theta
\]
就不应主要理解为数据在三个存储位置之间搬运。

更准确地说，这是认知结构的动力学相变。

从 $h$ 到 $E$：
\begin{statementbox}
\centering
\adjustbox{max width=.96\linewidth}{\(\displaystyle
\text{Activity}
\rightarrow
\text{EpisodeIdentity}.
\)}
\end{statementbox}

即当前状态退出实时活动，但获得更加稳定的历史身份。

从 $E$ 到 $\Theta$：
\begin{statementbox}
\centering
\adjustbox{max width=.96\linewidth}{\(\displaystyle
\text{EpisodeIdentity}
\rightarrow
\text{GenerativeStructure}.
\)}
\end{statementbox}

即具体事件个体性被压缩，稳定规律进入未来生成结构。

因此：
\begin{statementbox}
\centering
\adjustbox{max width=.96\linewidth}{\(\displaystyle
\text{MemoryTransfer}
\rightarrow
\text{MemoryPhaseTransition}.
\)}
\end{statementbox}

这也意味着，相变可以是渐进的，而不是瞬时完成的。

\section{遗忘不是单一现象}

三相理论还要求重新区分三种不同层次的“遗忘”。

首先，在 $h$ 中：
\[
h_C\rightarrow0
\]
通常只是一个结构退出当前活动。

这更准确地叫：
\begin{statementbox}
\centering
\adjustbox{max width=.96\linewidth}{\(\displaystyle
\text{Deactivation}.
\)}
\end{statementbox}

其次，在 $E$ 中：
\[
E_i\rightarrow0
\]
意味着某个具体事件不再能够被独立重新识别。

这可以称为：
\begin{statementbox}
\centering
\adjustbox{max width=.96\linewidth}{\(\displaystyle
\text{EpisodeDissolution}.
\)}
\end{statementbox}

第三，在 $\Theta$ 中：
\[
\Theta\rightarrow\Theta'
\]
意味着系统的生成结构、规律或技能本身发生变化。

这才属于：
\begin{statementbox}
\centering
\adjustbox{max width=.96\linewidth}{\(\displaystyle
\text{StructuralForgetting}.
\)}
\end{statementbox}

因此，当我们说一个智能体“忘记了”某件事情时，必须进一步询问：
\[
\text{是当前没有激活？}
\]
\[
\text{是无法回忆某次具体经历？}
\]
还是：
\[
\text{连从经历中学到的结构也已经消失？}
\]

三者对应完全不同的认知状态。

\section{三相记忆与认知结构对象}

为了进一步统一三相记忆，可以暂时不考虑某种具体内容，而设一个抽象认知结构为：
\[
C.
\]

同一个 $C$ 可以在不同时间以不同方式存在。

当它直接参与当前活动时：
\[
C@h.
\]

当它作为一次历史关系实例的一部分存在时：
\[
C@E.
\]

当它成为长期生成规则的一部分时：
\[
C@\Theta.
\]

因此，$h$、$E$、$\Theta$ 并不意味着系统分别存放三套完全不同的认知内容。

更加一般地：
\begin{statementbox}
\centering
\adjustbox{max width=.96\linewidth}{\(\displaystyle
\text{同一个认知结构可以以不同记忆相态存在。}
\)}
\end{statementbox}

例如一个对象概念可以当前被看到、曾经出现在某个具体事件中，同时又已经成为长期世界结构的一部分。

这意味着：
\begin{statementbox}
\centering
\adjustbox{max width=.96\linewidth}{\(\displaystyle
\text{RepresentationType}
\)}
\end{statementbox}
与：
\begin{statementbox}
\centering
\adjustbox{max width=.96\linewidth}{\(\displaystyle
\text{MemoryPhase}
\)}
\end{statementbox}
是两个正交维度。

对象、空间、关系、目标、自我等不同认知结构，都原则上可以分别存在于 $h$、$E$ 或 $\Theta$ 中。

\section{三相记忆首先是一种结构持续化理论}

由上述分析可以进一步得到：三相记忆真正研究的，并不仅仅是“记忆”。

它实际上研究：
\begin{statementbox}
\centering
\adjustbox{max width=.96\linewidth}{\(\displaystyle
\text{一个认知结构如何从瞬时出现，逐渐获得跨时间持续性，并最终成为系统自身的一部分。}
\)}
\end{statementbox}

这一过程可以写成：
\begin{statementbox}
\centering
\adjustbox{max width=.96\linewidth}{\(\displaystyle
\text{Emergence}
\rightarrow
\text{Persistence}
\rightarrow
\text{Structuralization}.
\)}
\end{statementbox}

对应：
\[
\boxed{
h
\rightarrow
E
\rightarrow
\Theta.
}
\]

因此，三相记忆可以被看成一种：
\begin{statementbox}
\centering
\adjustbox{max width=.96\linewidth}{\(\displaystyle
\text{CognitiveStructurePersistenceDynamics}.
\)}
\end{statementbox}

即“认知结构持续化动力学”。

从这一角度看，Memory 只是这种持续化动力学在认知系统中的主要表现。

\section{工作记忆有限性为什么会自然要求三相结构}

设工作记忆能够同时承载的有效认知结构复杂度存在上限：
\[
C_h(t)\le C_{\max}.
\]

若一个持续运行的 Agent 不断接收新经验：
\[
\frac{dI}{dt}>0,
\]
而旧结构始终保持在 $h$ 中，则长期必然出现：
\[
C_h(t)>C_{\max}.
\]

因此，一个有限工作记忆但需要持续运行的系统，必须存在：
\begin{statementbox}
\centering
\adjustbox{max width=.96\linewidth}{\(\displaystyle
h\rightarrow\text{non-}h
\)}
\end{statementbox}
机制。

然而，如果所有退出 $h$ 的结构都直接被删除，那么过去便无法继续影响未来，系统也无法形成真正的长期学习。

因此，至少需要存在：
\begin{statementbox}
\centering
\adjustbox{max width=.96\linewidth}{\(\displaystyle
\text{ActiveState}
\rightarrow
\text{PersistentNonActiveState}.
\)}
\end{statementbox}

这对应 $E$ 或其他较慢状态的必要性。

但如果所有历史经历都以完整 Episode 永久累积：
\[
|E(t)|\propto t,
\]
长期同样会产生无限增长问题。

因此系统进一步需要：
\begin{statementbox}
\centering
\adjustbox{max width=.96\linewidth}{\(\displaystyle
E
\rightarrow
\text{CompressedGenerativeStructure}.
\)}
\end{statementbox}

即：
\[
\boxed{
E\rightarrow\Theta.
}
\]

由此可以得到一个重要的理论推导：

如果一个系统同时满足：
\begin{statementbox}
\centering
\adjustbox{max width=.96\linewidth}{\(\displaystyle
\text{有限工作记忆}
\)}
\end{statementbox}
\begin{statementbox}
\centering
\adjustbox{max width=.96\linewidth}{\(\displaystyle
\text{持续生命周期}
\)}
\end{statementbox}
以及
\begin{statementbox}
\centering
\adjustbox{max width=.96\linewidth}{\(\displaystyle
\text{过去必须长期影响未来},
\)}
\end{statementbox}
那么它必须具有某种多时间尺度的结构沉降机制。

三相记忆
\[
h\rightarrow E\rightarrow\Theta
\]
正是这一机制的一种最小抽象形式。

\section{$E$ 是快速经历与慢结构之间的保护层}

上述推导还说明了另一个重要问题：为什么不能简单采用
\[
h\rightarrow\Theta.
\]

当前活动 $h$ 中包含大量：
\[
\text{噪声},
\quad
\text{偶然事件},
\quad
\text{单样本特例},
\quad
\text{临时目标},
\quad
\text{误解}.
\]

如果每一个当前状态都直接改写长期结构：
\[
h_t\rightarrow\Delta\Theta_t,
\]
系统将极度不稳定。

因此，$E$ 具有一个比“保存情景”更深的系统功能：
\begin{statementbox}
\centering
\adjustbox{max width=.96\linewidth}{\(\displaystyle
E=
\text{BufferBetweenExperienceAndSlowStructure}.
\)}
\end{statementbox}

即：
\begin{statementbox}
\centering
\adjustbox{max width=.96\linewidth}{\(\displaystyle
h
\rightarrow
E
\rightarrow
\text{RepeatedComparison}
\rightarrow
\text{Validation}
\rightarrow
\Theta.
\)}
\end{statementbox}

系统可以先保留一次经历，而暂时不把它当成普遍规律；以后通过更多经验重新评估其意义，最后才决定是否改变长期结构。

因此：
\begin{statementbox}
\centering
\adjustbox{max width=.96\linewidth}{\(\displaystyle
E
\text{ 是长期学习的历史缓冲与验证层。}
\)}
\end{statementbox}

这也是持续智能体能够同时获得可塑性与稳定性的关键机制之一。

\section{$\Theta$ 不是冻结结构，而是缓慢生成结构}

将 $\Theta$ 类比为“晶态”并不意味着：
\[
\Theta=\text{Frozen}.
\]

更准确地说：
\begin{statementbox}
\centering
\adjustbox{max width=.96\linewidth}{\(\displaystyle
\Theta=
\text{SlowlyReconfigurableGenerativeStructure}.
\)}
\end{statementbox}

即缓慢可重构的生成结构。

它与 $h$、$E$ 的主要差别是变化速度和因果地位：
\[
|\dot\Theta|
\ll
|\dot E|
\ll
|\dot h|.
\]

但长期生命周期中：
\[
\dot\Theta\neq0
\]
仍然可能成立。

因此，结构记忆既要足够稳定，以保持知识、技能和身份连续，又要保持有限可塑性，以适应长期环境变化。

三相记忆的关键不是：
\[
\text{变化与不变化},
\]
而是：
\begin{statementbox}
\centering
\adjustbox{max width=.96\linewidth}{\(\displaystyle
\text{不同速度的变化}.
\)}
\end{statementbox}

\section{$\Theta$ 为什么同时属于过去与未来}

如果 $\Theta$ 只是把过去压缩保存下来，那么它仍然只是另一种过去记忆。

但结构记忆真正重要的作用，是改变未来状态的概率、生成路径和可达空间。

因此可以抽象地写成：
\[
\boxed{
\Theta_t
\rightarrow
P(X_{t+\Delta}\mid X_t,\Theta_t).
}
\]

这意味着 $\Theta$ 将过去经验压缩成一个未来生成条件。

因此：
\begin{statementbox}
\centering
\adjustbox{max width=.96\linewidth}{\(\displaystyle
Past
\rightarrow
\Theta
\rightarrow
PossibleFuture.
\)}
\end{statementbox}

所以 $\Theta$ 可以被理解为：
\begin{statementbox}
\centering
\adjustbox{max width=.96\linewidth}{\(\displaystyle
\text{已经进入未来的过去}.
\)}
\end{statementbox}

这一点是三相记忆区别于传统记忆分类理论的核心之一。

记忆的最终功能并不是最大限度保存过去，而是：
\begin{statementbox}
\centering
\adjustbox{max width=.96\linewidth}{\(\displaystyle
\text{利用过去改变未来状态空间}.
\)}
\end{statementbox}

\section{三相记忆构成内部结构时间}

从三相角度看，智能体内部的时间不只是外部时钟的线性刻度。

对于智能体而言：
\[
h
\]
定义了当前活动面；

\[
E
\]
构成了可以重新进入现在的历史区域；

\[
\Theta
\]
则把过去压缩成未来生成约束。

因此可以写成：
\begin{statementbox}
\centering
\adjustbox{max width=.96\linewidth}{\(\displaystyle
E_{\mathrm{past}}
\rightarrow
h_{\mathrm{present}}
\rightarrow
\Theta_{\mathrm{future\ potential}}.
\)}
\end{statementbox}

更准确地说，$\Theta$ 不是未来本身，而是决定未来可能性分布的结构。

因此三相记忆共同构成一种内部结构时间：
\begin{statementbox}
\centering
\adjustbox{max width=.96\linewidth}{\(\displaystyle
\text{可识别过去}
+
\text{当前活动}
+
\text{未来生成潜势}.
\)}
\end{statementbox}

这也解释了为什么持续智能体的“时间”必须与记忆系统紧密相关：如果没有过去的结构保持，也没有未来的生成约束，那么系统只剩下一系列互不相干的瞬时状态。

\section{召回与巩固是相反方向的时间尺度迁移}

三相记忆的双向关系可以进一步统一为两种基本过程。

第一种是巩固：
\[
h\rightarrow E\rightarrow\Theta.
\]

其本质是：
\begin{statementbox}
\centering
\adjustbox{max width=.96\linewidth}{\(\displaystyle
\text{FastStructure}
\rightarrow
\text{SlowStructure}.
\)}
\end{statementbox}

即认知结构逐渐获得更长时间常数和更稳定的承载。

第二种是召回或激活：
\[
\Theta\rightarrow E\rightarrow h
\]
或直接：
\[
\Theta\rightarrow h,
\qquad
E\rightarrow h.
\]

其本质是：
\begin{statementbox}
\centering
\adjustbox{max width=.96\linewidth}{\(\displaystyle
\text{SlowStructure}
\rightarrow
\text{FastActiveStructure}.
\)}
\end{statementbox}

因此可以把记忆循环理解成：
\begin{statementbox}
\centering
\adjustbox{max width=.96\linewidth}{\(\displaystyle
\text{认知结构在不同时间尺度之间的双向迁移}.
\)}
\end{statementbox}

这比单纯“读取—写入”更能解释持续智能中的学习和回忆。

\section{$\Theta\rightarrow E$ 不只是回忆，还包括预测与想象}

如果 $\Theta$ 是生成结构，那么它不仅能够帮助恢复过去，也能够根据当前条件生成可能的事件轨迹。

可以写成：
\[
\boxed{
\Theta
+
Context
\rightarrow
\hat E.
}
\]

这里的 $\hat E$ 可以具有多种认识论状态。

过去真实发生过的：
\[
E^{\mathrm{obs}}.
\]

未来预测的：
\[
E^{\mathrm{pred}}.
\]

用于比较的反事实：
\[
E^{\mathrm{cf}}.
\]

因此 $\Theta\rightarrow E$ 包括：
\[
\text{回忆},
\quad
\text{模拟},
\quad
\text{预测},
\quad
\text{反事实构造}.
\]

这说明三相记忆并不是只处理 Past，而是天然连接 Past、Present 与 Possible Future。

\section{三相记忆是循环，而不是单向层级}

如果仅绘制：
\[
h\rightarrow E\rightarrow\Theta,
\]
很容易把三相误解成一个线性存储金字塔。

更完整的关系是：
\[
\boxed{
h
\leftrightarrow
E
\leftrightarrow
\Theta.
}
\]

其中：
\[
h\rightarrow E
\]
形成经历；

\[
E\rightarrow h
\]
实现情景召回；

\[
E\rightarrow\Theta
\]
形成结构学习；

\[
\Theta\rightarrow E
\]
生成或重新解释情景；

\[
\Theta\rightarrow h
\]
直接塑造当前认知；

在某些严格受控情况下：
\[
h\rightarrow\Theta
\]
还可能形成快速结构修正。

于是完整三相循环可以理解为：
\[
\boxed{
\Theta_t
\rightarrow
h_t
\rightarrow
E_t
\rightarrow
\Theta_{t+1}.
}
\]

也就是：
\begin{statementbox}
\centering
\adjustbox{max width=.96\linewidth}{\(\displaystyle
\text{已有结构产生当前经验，
当前经验形成历史，
历史重新塑造结构}.
\)}
\end{statementbox}

从这个角度看，三相记忆不是一个 Memory Stack，而是：
\begin{statementbox}
\centering
\adjustbox{max width=.96\linewidth}{\(\displaystyle
\text{持续主体内部结构的再生产循环}.
\)}
\end{statementbox}

\section{三相记忆的三个序参量}

为了以后对三相状态进行更严格的数学化，可以对任意认知结构 $C$ 定义三个概念性序参量。

第一，当前激活度：
\[
A(C,t),
\]
表示该结构是否直接参与当前认知。

第二，情景身份保持度：
\[
I_E(C,t),
\]
表示该结构是否仍然属于一个可以重新识别的具体历史事件。

第三，生成约束强度：
\[
G_\Theta(C,t),
\]
表示该结构对未来认知状态转移具有多强的稳定约束和生成能力。

于是典型 $h$ 相表现为：
\[
A(C,t)\gg0.
\]

典型 $E$ 相表现为：
\[
A(C,t)\approx0,
\qquad
I_E(C,t)\gg0.
\]

典型 $\Theta$ 相表现为：
\[
A(C,t)\approx0,
\qquad
G_\Theta(C,t)\gg0.
\]

这种定义说明三相并不只由时间长度决定，还涉及一个结构在系统中的功能地位。

因此，未来更完整的相态判定可以写成：
\begin{statementbox}
\centering
\adjustbox{max width=.96\linewidth}{\(\displaystyle
Phase(C,t)
=
f(
A,
I_E,
G_\Theta,
\tau
).
\)}
\end{statementbox}

\section{三相的最深层区别：现在、历史身份与生成规律}

综合上述分析，三相记忆可以压缩成三个最基本的问题。

对于 $h$：
\begin{statementbox}
\centering
\adjustbox{max width=.96\linewidth}{\(\displaystyle
\text{它是否仍然正在发生？}
\)}
\end{statementbox}

对于 $E$：
\begin{statementbox}
\centering
\adjustbox{max width=.96\linewidth}{\(\displaystyle
\text{我是否仍然能够识别“那一次发生了什么”？}
\)}
\end{statementbox}

对于 $\Theta$：
\begin{statementbox}
\centering
\adjustbox{max width=.96\linewidth}{\(\displaystyle
\text{它是否已经改变以后事情更可能怎样发生？}
\)}
\end{statementbox}

因此可以得到三相的简洁本体定义：
\begin{statementbox}
\centering
\adjustbox{max width=.96\linewidth}{\(\displaystyle
h:
\text{正在发生}.
\)}
\end{statementbox}
\begin{statementbox}
\centering
\adjustbox{max width=.96\linewidth}{\(\displaystyle
E:
\text{发生过}.
\)}
\end{statementbox}
\begin{statementbox}
\centering
\adjustbox{max width=.96\linewidth}{\(\displaystyle
\Theta:
\text{因为发生过，所以以后会不同}.
\)}
\end{statementbox}

这三个定义比“短期、中期、长期”更准确，也比“气态、液态、晶态”的比喻更加接近三相记忆真正的动力学含义。

\section{本章结论}

三相记忆的深度理解要求彻底摆脱“记忆等于存储”的传统直觉。

三相真正描述的是：
\begin{statementbox}
\centering
\adjustbox{max width=.96\linewidth}{\(\displaystyle
\text{过去在持续智能体中以何种动力深度继续存在}.
\)}
\end{statementbox}

其中：
\begin{statementbox}
\centering
\adjustbox{max width=.96\linewidth}{\(\displaystyle
h=
\text{仍直接参与当前因果闭环的认知结构},
\)}
\end{statementbox}
\begin{statementbox}
\centering
\adjustbox{max width=.96\linewidth}{\(\displaystyle
E=
\text{已经退出当前活动，但仍保持事件身份和可重激活性的历史结构},
\)}
\end{statementbox}
\begin{statementbox}
\centering
\adjustbox{max width=.96\linewidth}{\(\displaystyle
\Theta=
\text{由历史沉降形成，并已经成为未来生成条件的慢结构}.
\)}
\end{statementbox}

因此，三相记忆不是三个独立模块，而是同一认知结构持续化过程中的三个主要动力相区。

其核心过程不是：
\[
\text{保存}
\rightarrow
\text{搬运}
\rightarrow
\text{长期保存},
\]
而是：
\begin{statementbox}
\centering
\adjustbox{max width=.96\linewidth}{\(\displaystyle
\text{Activity}
\rightarrow
\text{EpisodeIdentity}
\rightarrow
\text{GenerativeStructure}.
\)}
\end{statementbox}

这一过程又可以反向进行：
\begin{statementbox}
\centering
\adjustbox{max width=.96\linewidth}{\(\displaystyle
\text{GenerativeStructure}
\rightarrow
\text{Episode}
\rightarrow
\text{CurrentActivity}.
\)}
\end{statementbox}

所以三相记忆最终是一套关于：
\begin{statementbox}
\centering
\adjustbox{max width=.96\linewidth}{\(\displaystyle
\text{认知结构如何出现、持续、沉降、重新激活并塑造未来}
\)}
\end{statementbox}
的多时间尺度动力学理论。

从这个角度看，可以给三相记忆一个更加基础的总结：
\begin{statementbox}
\centering
\adjustbox{max width=.96\linewidth}{\(\displaystyle
\text{\bfseries
记忆不是保存过去，而是让过去以不同速度继续存在于未来。
}
\)}
\end{statementbox}

而三相分别回答：
\begin{statementbox}
\centering
\adjustbox{max width=.96\linewidth}{\(\displaystyle
\text{\bfseries
$h$ 是仍然作为现在存在的过去；
$E$ 是仍然能够被识别的过去；
$\Theta$ 是已经成为未来生成结构的过去。
}
\)}
\end{statementbox}
